
\chapter{Quasiprimitive Groups and the O'Nan--Scott--Praeger Types}
\label{chap:quasiprimitive-ons-praeger}

\begin{topic}
For a finite quasiprimitive group \(G\leq \operatorname{Sym}(\Omega)\), the structure of
the minimal normal subgroups strongly restricts the whole permutation group.

The O'Nan--Scott--Praeger theorem divides such groups into eight types. The first
division is:

\begin{itemize}
    \item \(G\) has a unique minimal normal subgroup;
    \item \(G\) has exactly two minimal normal subgroups.
\end{itemize}

The second case leads to the holomorph simple and holomorph compound types.
The first case leads to affine, almost simple, twisted wreath, diagonal, and product
action types.
\end{topic}

\section{Case I: A Unique Minimal Normal Subgroup}

Assume first that \(G\leq \operatorname{Sym}(\Omega)\) is quasiprimitive and has a unique
minimal normal subgroup \(N\).

\subsection{The Centralizer of a Nonabelian Minimal Normal Subgroup}

\begin{lemma}
Let \(G\) have a unique minimal normal subgroup \(N\). If \(N\) is nonabelian, then
\[
    C_G(N)=1.
\]
Consequently, conjugation gives an embedding
\[
    G\hookrightarrow \operatorname{Aut}(N).
\]
\end{lemma}

\begin{proof}
Since \(N\unlhd G\), its centralizer \(C_G(N)\) is normal in \(G\).

Suppose \(C_G(N)\neq 1\). Since \(G\) is finite, \(C_G(N)\) contains a minimal normal
subgroup of \(G\). But \(N\) is the unique minimal normal subgroup of \(G\), so
\[
    N\leq C_G(N).
\]
Thus every element of \(N\) centralizes \(N\), meaning
\[
    N\leq Z(N).
\]
Hence \(N\) is abelian, contradicting the assumption that \(N\) is nonabelian.

Therefore
\[
    C_G(N)=1.
\]

Now \(G\) acts on \(N\) by conjugation. The kernel of this action is exactly \(C_G(N)\).
Hence
\[
    G/C_G(N)\hookrightarrow \operatorname{Aut}(N).
\]
Since \(C_G(N)=1\), we get
\[
    G\hookrightarrow \operatorname{Aut}(N).
\]
\end{proof}

\subsection{The Abelian Case: HA Type}

\begin{proposition}
Let \(G\leq \operatorname{Sym}(\Omega)\) be quasiprimitive, and suppose that \(G\) has a
unique minimal normal subgroup
\[
    N\cong C_p^d
\]
for some prime \(p\) and integer \(d\geq 1\). Then \(N\) is regular on \(\Omega\).
\end{proposition}

\begin{proof}
Since \(N\unlhd G\) and \(G\) is quasiprimitive, \(N\) is transitive on \(\Omega\).

Fix \(\omega\in\Omega\). Let \(n\in N_\omega\). For any \(\alpha\in\Omega\), since \(N\)
is transitive, there exists \(m\in N\) such that
\[
    \alpha=\omega^m.
\]
Because \(N\) is abelian, we have
\[
    \alpha^n
    =
    (\omega^m)^n
    =
    (\omega^n)^m
    =
    \omega^m
    =
    \alpha.
\]
Thus \(n\) fixes every point of \(\Omega\). Since the action is faithful, \(n=1\). Therefore
\[
    N_\omega=1.
\]
Hence \(N\) is regular.
\end{proof}

\begin{definition}[HA type]
A quasiprimitive group \(G\leq \operatorname{Sym}(\Omega)\) is said to be of
\emph{holomorph affine type}, or \(\mathrm{HA}\) type, if \(G\) has a unique minimal
normal subgroup
\[
    N\cong C_p^d
\]
and \(N\) is regular on \(\Omega\).

Equivalently, after identifying \(\Omega\) with the vector space \(N\cong \F_p^d\), the group
\(G\) embeds into the affine general linear group
\[
    \operatorname{AGL}(d,p)=\F_p^d\rtimes \operatorname{GL}(d,p).
\]
\end{definition}

\begin{example}[An HA-type quasiprimitive group]
    Let
    \[
        \Omega=\mathbb F_3
    \]
    and let
    \[
        G=\operatorname{AGL}(1,3)
        =
        \mathbb F_3^+ \rtimes \mathbb F_3^\times
    \]
    act on \(\Omega\) by affine transformations
    \[
        x\longmapsto ax+b,
        \qquad
        a\in \mathbb F_3^\times,\ b\in \mathbb F_3.
    \]
    Thus
    \[
        |G|=3\cdot 2=6,
    \]
    and in fact
    \[
        G\cong S_3.
    \]
    
    Let
    \[
        N=\{x\mapsto x+b\mid b\in\mathbb F_3\}.
    \]
    Then
    \[
        N\cong \mathbb F_3^+\cong C_3.
    \]
    The subgroup \(N\) is normal in \(G\), since translations are normalized by affine
    transformations. Moreover, \(N\) acts regularly on \(\Omega\): for any
    \(\alpha,\beta\in\mathbb F_3\), there is a unique translation
    \[
        x\mapsto x+(\beta-\alpha)
    \]
    sending \(\alpha\) to \(\beta\).
    
    Since \(G\cong S_3\), the subgroup corresponding to \(N\) is \(A_3\). Hence \(N\) is the
    unique minimal normal subgroup of \(G\). Therefore every nontrivial normal subgroup of
    \(G\) contains \(N\), and is transitive on \(\Omega\). Thus \(G\) is quasiprimitive.
    
    Consequently, \(G=\operatorname{AGL}(1,3)\) is of \(\mathrm{HA}\) type, with regular
    abelian minimal normal subgroup
    \[
        N\cong C_3.
    \]
    \end{example}

\subsection{The Nonabelian Simple Case: AS Type}

\begin{proposition}
Let \(G\leq \operatorname{Sym}(\Omega)\) be quasiprimitive, and suppose that \(G\) has a
unique minimal normal subgroup \(N\), where \(N\) is nonabelian simple. Then
\[
    N\unlhd G\leq \operatorname{Aut}(N).
\]
\end{proposition}

\begin{proof}
Since \(N\) is nonabelian and is the unique minimal normal subgroup of \(G\), the previous
lemma gives
\[
    C_G(N)=1.
\]
Therefore the conjugation action of \(G\) on \(N\) embeds \(G\) into \(\operatorname{Aut}(N)\).
Under this embedding, \(N\) is identified with the inner automorphism group
\[
    \operatorname{Inn}(N)\cong N.
\]
Thus
\[
    N\unlhd G\leq \operatorname{Aut}(N).
\]
\end{proof}

\begin{definition}[AS type]
A quasiprimitive group \(G\leq \operatorname{Sym}(\Omega)\) is said to be of
\emph{almost simple type}, or \(\mathrm{AS}\) type, if \(G\) has a unique minimal normal
subgroup \(N\), where \(N\) is nonabelian simple, and
\[
    N\unlhd G\leq \operatorname{Aut}(N).
\]
\end{definition}

\subsection{The Nonabelian Characteristically Simple Case}

Now suppose
\[
    N=T^k,\qquad k\geq 2,
\]
where \(T\) is nonabelian simple. Let \(\omega\in\Omega\). Since \(N\unlhd G\) and
\(G\) is quasiprimitive, \(N\) is transitive. Thus \(\Omega\) may be identified with the set of
right cosets of \(N_\omega\) in \(N\).

The possible structure of \(N_\omega\) leads to four types.

\begin{definition}[Full diagonal subgroup]
Let
\[
    N=T_1\times T_2\times \cdots \times T_k,
\]
where each \(T_i\cong T\). A subgroup \(D\leq N\) is called a \emph{full diagonal subgroup}
over a subset \(I\subseteq \{1,\dots,k\}\) if, after choosing isomorphisms
\[
    \alpha_i:T\to T_i\qquad (i\in I),
\]
we have
\[
    D=
    \{(t^{\alpha_i})_{i\in I}\mid t\in T\}
    \leq \prod_{i\in I}T_i.
\]
Thus \(D\cong T\) and projects onto each direct factor \(T_i\) for \(i\in I\).
\end{definition}

\begin{topic}
Let \(N=T^k\), \(k\geq 2\), be the unique minimal normal subgroup of a quasiprimitive
group \(G\leq \operatorname{Sym}(\Omega)\). For \(\omega\in\Omega\), the point stabilizer
\(N_\omega\) determines the following four types:

\begin{enumerate}[label=(\roman*)]
    \item If
    \[
        N_\omega=1,
    \]
    then \(N\) is regular. This is the \(\mathrm{TW}\) type.

    \item If
    \[
        N_\omega\cong T
    \]
    is a full diagonal subgroup of \(T^k\), then this is the \(\mathrm{SD}\) type.

    \item If
    \[
        N_\omega\cong T^m,\qquad 2\leq m\leq \frac{k}{2},
    \]
    and \(N_\omega\) is a product of full diagonal subgroups corresponding to a partition
    of the \(k\) simple direct factors, then this is the \(\mathrm{CD}\) type.

    \item If
    \[
        1\neq N_\omega
    \]
    and \(N_\omega\) is not of the diagonal forms above, then this is the \(\mathrm{PA}\) type.
\end{enumerate}
\end{topic}

\begin{definition}[TW type]
A quasiprimitive group \(G\leq \operatorname{Sym}(\Omega)\) is said to be of
\emph{twisted wreath type}, or \(\mathrm{TW}\) type, if its unique minimal normal subgroup is
\[
    N=T^k,\qquad k\geq 2,
\]
where \(T\) is nonabelian simple, and \(N\) is regular on \(\Omega\).
\end{definition}

\begin{definition}[SD type]
A quasiprimitive group \(G\leq \operatorname{Sym}(\Omega)\) is said to be of
\emph{simple diagonal type}, or \(\mathrm{SD}\) type, if its unique minimal normal subgroup is
\[
    N=T^k,\qquad k\geq 2,
\]
where \(T\) is nonabelian simple, and for \(\omega\in\Omega\),
\[
    N_\omega\cong T
\]
is a full diagonal subgroup of \(N\).
\end{definition}

\begin{definition}[CD type]
A quasiprimitive group \(G\leq \operatorname{Sym}(\Omega)\) is said to be of
\emph{compound diagonal type}, or \(\mathrm{CD}\) type, if its unique minimal normal subgroup is
\[
    N=T^k,\qquad k\geq 2,
\]
where \(T\) is nonabelian simple, and for \(\omega\in\Omega\),
\[
    N_\omega\cong T^m,\qquad 2\leq m\leq \frac{k}{2},
\]
where \(N_\omega\) is a product of full diagonal subgroups.
\end{definition}

\begin{definition}[PA type]
A quasiprimitive group \(G\leq \operatorname{Sym}(\Omega)\) is said to be of
\emph{product action type}, or \(\mathrm{PA}\) type, if its unique minimal normal subgroup is
\[
    N=T^k,\qquad k\geq 2,
\]
where \(T\) is nonabelian simple, and \(N_\omega\neq 1\) is not a diagonal stabilizer of
the \(\mathrm{SD}\) or \(\mathrm{CD}\) form.
\end{definition}

\section{Regular Representations and Central Products}

This section collects the regular-action machinery used in the two-minimal-normal-subgroup
case.

\subsection{Right Regular, Left Regular, and Conjugation Actions}

Let \(G\) be a finite group. Define three families of permutations of the set \(G\).

For \(g\in G\), define
\[
    x^{\widehat g}=xg,
\]
\[
    x^{\check g}=g^{-1}x,
\]
and
\[
    x^{\overline g}=g^{-1}xg.
\]
Set
\[
    \widehat G=\{\widehat g\mid g\in G\},
\]
\[
    \check G=\{\check g\mid g\in G\},
\]
and
\[
    \overline G=\{\overline g\mid g\in G\}.
\]

\begin{lemma}[Right and left regular representations]
Let \(G\) be a finite group. Then the following hold.

\begin{enumerate}[label=(\arabic*)]
    \item \(\widehat G,\check G\leq \operatorname{Sym}(G)\), and both are regular.

    \item \(\widehat G\) and \(\check G\) centralize each other:
    \[
        \widehat g\check h=\check h\widehat g
        \qquad (g,h\in G).
    \]

    \item
    \[
        \widehat G\cap \check G
        =
        Z(\widehat G)
        =
        Z(\check G)
        =
        \widehat{Z(G)}
        =
        \check{Z(G)}.
    \]

    \item If
    \[
        X=\langle \widehat G,\check G\rangle,
    \]
    then
    \[
        X=\widehat G\check G
    \]
    and
    \[
        X\cong
        (\widehat G\times \check G)/
        \{(\widehat z,\check z)\mid z\in Z(G)\}.
    \]

    \item If \(1\in G\) is used as base point, then
    \[
        X_1=\overline G.
    \]
\end{enumerate}
\end{lemma}

\begin{proof}
For \(g,h\in G\),
\[
    x^{\widehat g\widehat h}=(xg)h=x(gh)=x^{\widehat{gh}},
\]
so
\[
    \widehat g\widehat h=\widehat{gh}.
\]
Hence \(\widehat G\leq \operatorname{Sym}(G)\). Similarly,
\[
    x^{\check g\check h}
    =
    h^{-1}(g^{-1}x)
    =
    (gh)^{-1}x
    =
    x^{\check{gh}},
\]
so
\[
    \check g\check h=\check{gh}.
\]
Thus \(\check G\leq \operatorname{Sym}(G)\).

The group \(\widehat G\) is regular because for any \(x,y\in G\), there is a unique \(g\in G\)
such that
\[
    xg=y,
\]
namely \(g=x^{-1}y\). Similarly, \(\check G\) is regular.

Next,
\[
    x^{\widehat g\check h}
    =
    h^{-1}xg
    =
    x^{\check h\widehat g}.
\]
Therefore
\[
    \widehat g\check h=\check h\widehat g.
\]

Now suppose
\[
    \widehat g=\check h.
\]
Evaluating at \(1\in G\), we get
\[
    g=h^{-1}.
\]
Thus for all \(x\in G\),
\[
    xg=h^{-1}x=gx.
\]
Hence \(g\in Z(G)\). Conversely, if \(z\in Z(G)\), then
\[
    x^{\widehat z}=xz=zx=x^{\check{z^{-1}}},
\]
so \(\widehat z=\check{z^{-1}}\). Therefore
\[
    \widehat G\cap \check G=\widehat{Z(G)}=\check{Z(G)}.
\]
Also
\[
    Z(\widehat G)=\widehat{Z(G)},\qquad
    Z(\check G)=\check{Z(G)}.
\]

Since \(\widehat G\) and \(\check G\) centralize each other,
\[
    X=\langle \widehat G,\check G\rangle=\widehat G\check G.
\]
Consider the epimorphism
\[
    \theta:\widehat G\times \check G\to X,
    \qquad
    (\widehat g,\check h)\mapsto \widehat g\check h.
\]
Its kernel consists of those pairs satisfying
\[
    \widehat g\check h=1.
\]
For all \(x\in G\),
\[
    x^{\widehat g\check h}=h^{-1}xg.
\]
Thus \(\widehat g\check h=1\) if and only if
\[
    h^{-1}xg=x
    \qquad \forall x\in G.
\]
Taking \(x=1\), we get \(g=h\). Then
\[
    g^{-1}xg=x
    \qquad \forall x\in G,
\]
so \(g\in Z(G)\). Hence
\[
    \ker\theta=
    \{(\widehat z,\check z)\mid z\in Z(G)\}.
\]
The isomorphism theorem gives
\[
    X\cong
    (\widehat G\times \check G)/
    \{(\widehat z,\check z)\mid z\in Z(G)\}.
\]

Finally, an element of \(X\) has the form \(\widehat g\check h\). It fixes \(1\in G\) if and
only if
\[
    1^{\widehat g\check h}=h^{-1}g=1,
\]
that is, if and only if \(g=h\). Therefore
\[
    X_1=\{\widehat g\check g\mid g\in G\}.
\]
But
\[
    x^{\widehat g\check g}
    =
    g^{-1}xg
    =
    x^{\overline g}.
\]
Hence
\[
    X_1=\overline G.
\]
\end{proof}

\begin{corollary}
If \(Z(G)=1\), then
\[
    \widehat G\cap \check G=1
\]
and hence
\[
    \langle \widehat G,\check G\rangle
    =
    \widehat G\times \check G.
\]
Moreover,
\[
    \overline G\cong G.
\]
\end{corollary}

\subsection{Central Products}

\begin{definition}[Central product]
Let \(M\) and \(N\) be groups, and let
\[
    C_1\leq Z(M),\qquad C_2\leq Z(N).
\]
Suppose that
\[
    \varphi:C_1\to C_2
\]
is an isomorphism. Define
\[
    D_\varphi=
    \{(c,\varphi(c)^{-1})\mid c\in C_1\}
    \leq M\times N.
\]
The quotient
\[
    (M\times N)/D_\varphi
\]
is called the \emph{central product} of \(M\) and \(N\) amalgamating \(C_1\) and \(C_2\).

When \(C_1\cong C_2\cong C\), we often write this central product as
\[
    M\circ_C N
\]
or simply
\[
    M\circ N
\]
when the amalgamated central subgroups are clear.
\end{definition}

\begin{remark}
A central product is obtained from the direct product \(M\times N\) by identifying selected
central subgroups of \(M\) and \(N\). Thus it is generally not a direct product, but it is still
built from two commuting subgroups.
\end{remark}

\begin{example}
Let \(Q_8=\{\pm1,\pm i,\pm j,\pm k\}\). Then
\[
    Z(Q_8)=\{1,-1\}\cong C_2.
\]
Taking two copies of \(Q_8\), identify the two central involutions. Then
\[
    Q_8\circ Q_8
    =
    (Q_8\times Q_8)/\langle(-1,-1)\rangle.
\]
Hence
\[
    |Q_8\circ Q_8|
    =
    \frac{|Q_8|^2}{2}
    =
    \frac{8\cdot 8}{2}
    =
    32.
\]
\end{example}

\begin{example}
Let
\[
    C_4=\gen{c}.
\]
Since
\[
    \gen{c^2}\leq Z(C_4),
\]
we may identify
\[
    \gen{-1}\leq Z(Q_8)
\]
with
\[
    \gen{c^2}\leq Z(C_4).
\]
Then
\[
    Q_8\circ C_4
    =
    (Q_8\times C_4)/\langle(-1,c^2)\rangle.
\]
\end{example}

\begin{example}
Let
\[
    M=\gen{a,z\mid a^8=z^2=1,\ a^z=a^5},
\]
where
\[
    a^z=z^{-1}az.
\]
Then
\[
    (a^2)^z=(a^z)^2=(a^5)^2=a^{10}=a^2,
\]
so
\[
    \gen{a^2}\leq Z(M).
\]
In fact,
\[
    Z(M)=\gen{a^2}\cong C_4.
\]

Similarly, let
\[
    N=\gen{b,y\mid b^8=y^2=1,\ b^y=b^5}.
\]
Then
\[
    Z(N)=\gen{b^2}\cong C_4.
\]

If we identify
\[
    \gen{a^2}\leq Z(M)
\]
with
\[
    \gen{b^2}\leq Z(N),
\]
then we obtain the central product
\[
    M\circ N
    =
    (M\times N)/\langle(a^2,b^{-2})\rangle.
\]

If instead we identify only the central subgroups of order \(2\),
\[
    \gen{a^4}\leq Z(M),
    \qquad
    \gen{b^4}\leq Z(N),
\]
then
\[
    M\circ_{C_2}N
    =
    (M\times N)/\langle(a^4,b^4)\rangle.
\]
\end{example}

\section{Centralizers of Regular Subgroups}

\begin{lemma}
Let
\[
    M\leq \operatorname{Sym}(\Omega)
\]
be regular, and fix \(\omega_0\in\Omega\). Identify \(\Omega\) with \(M\) by
\[
    \omega_0^m\longmapsto m.
\]
Under this identification, the action of \(M\) on \(\Omega\) becomes the right regular action
\[
    \widehat M=\{\widehat m\mid m\in M\}
\]
on the set \(M\).

Moreover,
\[
    C_{\operatorname{Sym}(\Omega)}(M)
    \cong
    C_{\operatorname{Sym}(M)}(\widehat M)
    =
    \check M.
\]
In particular, \(C_{\operatorname{Sym}(\Omega)}(M)\) is regular.
\end{lemma}

\begin{proof}
Since \(M\) is regular, for each \(\omega\in\Omega\) there exists a unique \(m\in M\) such that
\[
    \omega=\omega_0^m.
\]
Thus
\[
    \omega_0^m\longmapsto m
\]
is a bijection \(\Omega\to M\).

For \(g\in M\),
\[
    (\omega_0^m)^g=\omega_0^{mg}.
\]
Under the identification \(\Omega\cong M\), this is
\[
    m\longmapsto mg,
\]
which is the right regular action \(\widehat g\).

Now let
\[
    \alpha\in C_{\operatorname{Sym}(M)}(\widehat M).
\]
Set
\[
    1^\alpha=t.
\]
Then for every \(m\in M\),
\[
    m^\alpha
    =
    (1^{\widehat m})^\alpha
    =
    (1^\alpha)^{\widehat m}
    =
    tm.
\]
Thus \(\alpha\) is left multiplication by \(t\), which is \(\check{t^{-1}}\). Hence
\[
    \alpha\in \check M.
\]
Conversely, every left multiplication map commutes with every right multiplication map.
Therefore
\[
    C_{\operatorname{Sym}(M)}(\widehat M)=\check M.
\]
\end{proof}

\section{Case II: Two Minimal Normal Subgroups}

Now suppose \(G\leq \operatorname{Sym}(\Omega)\) is quasiprimitive and has two distinct
minimal normal subgroups \(M\) and \(N\).

\begin{proposition}
Let \(G\leq \operatorname{Sym}(\Omega)\) be quasiprimitive. If \(M\) and \(N\) are two distinct
minimal normal subgroups of \(G\), then:

\begin{enumerate}[label=(\arabic*)]
    \item
    \[
        M\cap N=1;
    \]

    \item
    \[
        [M,N]=1;
    \]

    \item \(M\) and \(N\) are regular on \(\Omega\);

    \item
    \[
        M\cong N;
    \]

    \item There exists a nonabelian simple group \(T\) and an integer \(k\geq 1\) such that
    \[
        M\cong N\cong T^k.
    \]
\end{enumerate}
\end{proposition}

\begin{proof}
Since \(M,N\unlhd G\), their intersection is normal in \(G\):
\[
    M\cap N\unlhd G.
\]
Moreover,
\[
    M\cap N\leq M
    \qquad\text{and}\qquad
    M\cap N\leq N.
\]
By minimal normality and the assumption \(M\neq N\), we get
\[
    M\cap N=1.
\]

Now
\[
    [M,N]\leq M\cap N,
\]
so
\[
    [M,N]=1.
\]
Thus \(M\) and \(N\) centralize each other.

Since \(G\) is quasiprimitive and \(M,N\) are nontrivial normal subgroups, both \(M\) and
\(N\) are transitive on \(\Omega\).

Fix \(\omega\in\Omega\). Let \(n\in N_\omega\). For any \(\alpha\in\Omega\), because \(M\)
is transitive, there exists \(m\in M\) such that
\[
    \alpha=\omega^m.
\]
Since \(M\) and \(N\) centralize each other,
\[
    \alpha^n
    =
    (\omega^m)^n
    =
    (\omega^n)^m
    =
    \omega^m
    =
    \alpha.
\]
Thus \(n\) fixes every point of \(\Omega\). Since the action is faithful, \(n=1\). Therefore
\[
    N_\omega=1.
\]
Hence \(N\) is regular. The same argument shows that \(M\) is regular.

Since both are regular,
\[
    |M|=|\Omega|=|N|.
\]

Because \(M\) is regular,
\[
    C_{\operatorname{Sym}(\Omega)}(M)
\]
is also regular and is isomorphic to \(M\). Since \(N\) centralizes \(M\), we have
\[
    N\leq C_{\operatorname{Sym}(\Omega)}(M).
\]
Both groups have order \(|\Omega|\), so
\[
    N=C_{\operatorname{Sym}(\Omega)}(M).
\]
Therefore
\[
    N\cong M.
\]

Finally, minimal normal subgroups are characteristically simple, so
\[
    M\cong T^k
\]
for some finite simple group \(T\). If \(T\) were abelian, then \(M\) would be abelian.
Since \(M\) is regular and abelian, its centralizer in \(\operatorname{Sym}(\Omega)\) would be
itself. Hence
\[
    C_{\operatorname{Sym}(\Omega)}(M)=M.
\]
But we have already shown that
\[
    C_{\operatorname{Sym}(\Omega)}(M)=N,
\]
so \(M=N\), contradicting the assumption that \(M\) and \(N\) are distinct.

Thus \(T\) is nonabelian.
\end{proof}

\subsection{Relabelling \texorpdfstring{$\Omega$}{Omega} by a Regular Normal Subgroup}

Let \(M\leq \operatorname{Sym}(\Omega)\) be regular and fix
\[
    \Omega=\{\omega_1,\omega_2,\dots,\omega_n\}.
\]
Choose \(\omega_1\) as base point. Since \(M\) is regular, for each \(i\) there exists a unique
\(x_i\in M\) such that
\[
    \omega_i=\omega_1^{x_i}.
\]
Thus
\[
    M=\{1=x_1,x_2,\dots,x_n\}.
\]
The map
\[
    \omega_i\longmapsto x_i
\]
identifies \(\Omega\) with \(M\).

Under this identification, \(M\) acts on itself by right multiplication:
\[
    x\mapsto xg.
\]
Hence
\[
    M=\widehat M.
\]

If \(N\) is the second minimal normal subgroup, then
\[
    N=C_{\operatorname{Sym}(\Omega)}(M)=\check M.
\]
Thus
\[
    MN=\widehat M\check M.
\]

\begin{proposition}
With the above identification, let
\[
    X=MN=\widehat M\check M.
\]
Taking \(1\in M\) as base point, we have
\[
    X_1=\overline M
    =
    \{\widehat m\check m\mid m\in M\}.
\]
If \(M\cong T^k\) for a nonabelian simple group \(T\), then \(Z(M)=1\) and
\[
    X_1\cong M.
\]
\end{proposition}

\begin{proof}
Every element of \(X\) has the form
\[
    \widehat a\check b.
\]
Now
\[
    1^{\widehat a\check b}=b^{-1}a.
\]
Therefore
\[
    \widehat a\check b\in X_1
    \Longleftrightarrow
    b^{-1}a=1
    \Longleftrightarrow
    a=b.
\]
Hence
\[
    X_1=\{\widehat m\check m\mid m\in M\}.
\]
But
\[
    x^{\widehat m\check m}
    =
    m^{-1}xm
    =
    x^{\overline m}.
\]
Thus
\[
    X_1=\overline M.
\]

If \(M\cong T^k\) with \(T\) nonabelian simple, then \(Z(T)=1\), so
\[
    Z(M)=1.
\]
The conjugation action of \(M\) on itself has kernel \(Z(M)\), hence is faithful. Therefore
\[
    \overline M\cong M.
\]
\end{proof}

\section{The O'Nan--Scott--Praeger Theorem}

\begin{theorem}[O'Nan--Scott--Praeger theorem]
Let \(G\leq \operatorname{Sym}(\Omega)\) be a finite quasiprimitive permutation group.
Then \(G\) belongs to one of the following eight types.

\begin{description}[leftmargin=2.8cm, style=nextline]
    \item[\(\mathrm{HS}\)]
    \textbf{Holomorph Simple type.}

    The group \(G\) has exactly two minimal normal subgroups \(M\) and \(N\), and
    \[
        M\cong N\cong T,
    \]
    where \(T\) is a nonabelian simple group. Both \(M\) and \(N\) are regular.

    \item[\(\mathrm{HC}\)]
    \textbf{Holomorph Compound type.}

    The group \(G\) has exactly two minimal normal subgroups \(M\) and \(N\), and
    \[
        M\cong N\cong T^k,
        \qquad k\geq 2,
    \]
    where \(T\) is a nonabelian simple group. Both \(M\) and \(N\) are regular.

    \item[\(\mathrm{HA}\)]
    \textbf{Holomorph Affine type.}

    The group \(G\) has a unique minimal normal subgroup
    \[
        N\cong C_p^d
    \]
    for some prime \(p\), and \(N\) is regular.

    \item[\(\mathrm{AS}\)]
    \textbf{Almost Simple type.}

    The group \(G\) has a unique minimal normal subgroup \(N\), where \(N\) is nonabelian
    simple, and
    \[
        N\unlhd G\leq \operatorname{Aut}(N).
    \]

    \item[\(\mathrm{TW}\)]
    \textbf{Twisted Wreath type.}

    The group \(G\) has a unique minimal normal subgroup
    \[
        N=T^k,\qquad k\geq 2,
    \]
    where \(T\) is nonabelian simple, and \(N\) is regular.

    \item[\(\mathrm{SD}\)]
    \textbf{Simple Diagonal type.}

    The group \(G\) has a unique minimal normal subgroup
    \[
        N=T^k,\qquad k\geq 2,
    \]
    where \(T\) is nonabelian simple, and for \(\omega\in\Omega\),
    \[
        N_\omega\cong T
    \]
    is a full diagonal subgroup of \(N\).

    \item[\(\mathrm{CD}\)]
    \textbf{Compound Diagonal type.}

    The group \(G\) has a unique minimal normal subgroup
    \[
        N=T^k,\qquad k\geq 2,
    \]
    where \(T\) is nonabelian simple, and for \(\omega\in\Omega\),
    \[
        N_\omega\cong T^m,
        \qquad
        2\leq m\leq \frac{k}{2},
    \]
    where \(N_\omega\) is a product of full diagonal subgroups.

    \item[\(\mathrm{PA}\)]
    \textbf{Product Action type.}

    The group \(G\) has a unique minimal normal subgroup
    \[
        N=T^k,\qquad k\geq 2,
    \]
    where \(T\) is nonabelian simple, and \(N_\omega\neq 1\) is not of the diagonal forms
    appearing in \(\mathrm{SD}\) or \(\mathrm{CD}\).
\end{description}
\end{theorem}

\section{Summary of the Two-Minimal-Normal-Subgroup Case}

\begin{topic}
Suppose \(G\leq \operatorname{Sym}(\Omega)\) is quasiprimitive and has two distinct minimal
normal subgroups \(M\) and \(N\). Then:

\begin{enumerate}[label=(\arabic*)]
    \item
    \[
        M\cap N=1;
    \]

    \item
    \[
        [M,N]=1;
    \]

    \item \(M\) and \(N\) are regular;

    \item
    \[
        M\cong N\cong T^k
    \]
    for some nonabelian simple group \(T\);

    \item after identifying \(\Omega\) with \(M\),
    \[
        M=\widehat M,
        \qquad
        N=\check M;
    \]

    \item
    \[
        MN=\widehat M\check M;
    \]

    \item since \(Z(M)=1\),
    \[
        MN=\widehat M\times \check M;
    \]

    \item the stabilizer of \(1\in M\) in \(MN\) is
    \[
        (MN)_1=\overline M\cong M.
    \]
\end{enumerate}

If \(k=1\), this is the \(\mathrm{HS}\) type. If \(k\geq 2\), this is the \(\mathrm{HC}\) type.
\end{topic}